\subsection{Admissibility does not force a rational comparison map to have integral image}
\label{arch136:rational-image}

\paragraph{Archive and exact scope.}
This concerns the generalized divided-power de Rham theories of
cote136, physical pp.~117,119,121, archive footers116,118,120,
\texttt{ARCHIVE\_136\_scan117.tex},
\texttt{ARCHIVE\_136\_scan119.tex}, and
\texttt{ARCHIVE\_136\_scan121.tex}.
All input tags are none apart from the inherited admissibility of
the standard theory \(B\), AQ-136-DP-MODEL, and the elementary
Dold--Kan equivalence. There is no outside preprint input.

\paragraph{ARCH-136-RATIONAL-IMAGE-NOINTEGRAL.}
There is an admissible strict differential PD theory \(C\) over
\(S=\Gamma_{\mathbb Z}(T)\), with \(C_0=S\) and all underlying
abelian groups free, and a strict coefficient-preserving map
\[
                         \phi:C\longrightarrow B_{\mathbb Q}
\]
whose image is not contained in \(B\) in any weight tail.
Indeed one may prescribe an arbitrary integer denominator
\(d_m\geq2\) independently in every weight \(m\geq2\).
There are also strict maps
\[
                 B\xrightarrow{\iota}C\xrightarrow{\pi}B,
                    \qquad \pi\iota=\mathrm{id},
                    \qquad \phi\iota:B\hookrightarrow B_{\mathbb Q}.
\]
The map \(\phi\) is joined to the integral comparison
\(C\xrightarrow{\pi}B\hookrightarrow B_{\mathbb Q}\)
by a coefficient-preserving simplicial PD homotopy.
Thus the failure concerns the image of a specified rational map;
it is neither an obstruction to an integral map for this source
nor an obstruction to connectedness.

\paragraph{Construction and proof.}
Use normalization with differential \(d_0\) on the intersection
of the kernels of the other faces. Let \(U\) be the free simplicial
abelian group corresponding by Dold--Kan to the chain complex
\[
                \mathbb Z v[2]\xrightarrow{\ 1\ }
                \mathbb Z u[1].                         \tag{RI.1}
\]
It is simplicially contractible, \(U_0=0\), and each \(U_n\)
is finite free. For every \(m\geq2\), take a copy \(U^{(m)}\)
and give it weight \(m\). Define
\[
 \begin{split}
  V&=\bigoplus_{m\geq2}U^{(m)},\\
  F(V_n)&=\Gamma_{\mathbb Z}(a(V_n))
                      \otimes_{\mathbb Z}\Lambda_{\mathbb Z}(e(V_n)),
       \qquad |a|_{\rm form}=0,\quad |e|_{\rm form}=1,\\
  da(v)&=e(v),\qquad de(v)=0,\\
  C_n&=B_n\otimes_{\mathbb Z}F(V_n).
 \end{split}                                               \tag{RI.2}
\]
The coefficient algebra acts through \(B_n\), and the
divided powers are those of the free differential PD algebra.
Equivalently, the last tensor product is over \(S\) after
adjoining \(S\) to the second factor. All weights are nonnegative.
In each fixed weight there are only finitely many generators
contributing in any one simplex degree, so the displayed groups
are free abelian. The maps \(\iota\) and \(\pi\) are inclusion
and augmentation of the second factor.

Each even PD disk has the weight-preserving contraction
\[
 d(a^{[r]})=a^{[r-1]}e\quad(r\geq1),
\]
with coefficient one. Tensoring these actual contractions, and
then using the filtered union over finite sets of generators,
shows that \(C_n\to B_n\) is a vertical chain-homotopy
equivalence (weightwise, only finitely many generators matter).
Consequently its vertical cohomology is exactly \(S\), in
form degree zero, in every weight and every simplex degree.

The contraction of (RI.1), and hence of \(V\), is simplicial
and linear. Applying the functor \(F\) to the component maps
of its simplicial homotopy gives a simplicial strict differential
PD homotopy from \(\mathrm{id}_{F(V)}\) to its augmentation
followed by its unit. On the first tensor factor use the
appropriate degeneracy of \(B\). This gives a simplicial
strict PD homotopy
\[
                       \mathrm{id}_C\simeq\iota\pi .       \tag{RI.3}
\]
Every component preserves both gradings. Thus each row of
\(C\to B\) is a simplicial homotopy equivalence, with the same
horizontal vanishing bounds as \(B\). This proves the stated
admissibility, uniformly in simplex degree. Also \(C_0=B_0=S\).

Now fix the denominators \(d_m\). In the normalized zero-form
row of \(B_{\mathbb Q}\), the assignments
\[
 v\longmapsto \frac{x_1^{[m-1]}x_2}{d_m}
                   \in N_2(B_{\mathbb Q}^{m,0}),
 \qquad
 u\longmapsto \frac{x_0^{[m-1]}x_1}{d_m}
                   \in N_1(B_{\mathbb Q}^{m,0})             \tag{RI.4}
\]
define a chain map from (RI.1). Indeed the first expression
vanishes on faces1 and2, its face0 is the second expression,
and both faces of the second expression are zero. Dold--Kan
therefore gives a simplicial linear map
\(U^{(m)}\to B_{\mathbb Q}^{m,0}\).
Their direct sum extends, by the free PD universal property,
to (RI.2): send each \(a(v)\) to (RI.4), each \(e(v)\) to
its vertical differential, and use the usual inclusion on \(B\).
Since the target is rational its divided powers are the
canonical ones, and the special differential identity is
preserved. This constructs the claimed \(\phi\).

The monomial \(x_0^{[m-1]}x_1\) is a primitive basis vector
in \(B_1^{m,0}=\Gamma_{\mathbb Z}(x_0,x_1)^m\).
Its quotient by \(d_m\geq2\) is therefore not integral.
These nonintegral images occur in every weight \(m\geq2\),
proving the absence of any eventual image-integrality bound.
Finally, compose (RI.3) with \(\phi\). Its other endpoint is
\(\phi\iota\pi\), precisely the integral comparison asserted
in the statement. This proves the homotopy assertion.
\(\square\)

\paragraph{Why integral relative periods do not rule this out.}
The new zero-forms \(a(v)\) are not vertical cycles, and their
images in (RI.4) need not be integral polynomials. Their
differentials are exact forms. Thus the natural integral-period
comparison for relative cohomology remains valid: it tests
closed relative expressions, rather than arbitrary coefficient
vectors of a specified rational comparison map. Even a source
with standard complementary forms, free underlying groups,
and a retraction to \(B\) can carry such nonintegral maps.
An existence theorem for \emph{some} integral comparison map
would require an additional argument; the assertion that every
rational comparison already has integral image cannot supply it.

\paragraph{Verification status.}
The author has read the complete saved proof, including the
explicit normalized faces, both homotopies, and the uniform
admissibility verification. The parent and reviewer123 have each
independently read the complete saved proof and communicated PASS
on 2026-11-22. No main connectedness result is asserted by this
counterexample.
